% Formato del título de capítulos y secciones
% \titleformat{\chapter}[block]{\titlerule[0.8pt]\normalfont\huge\bfseries}{\thechapter.}{.5em}{\Huge}[{\titlerule[0.8pt]}]
% \titlespacing*{\chapter}{0pt}{-19pt}{25pt}
% \titleformat{\section}[block]{\normalfont\Large\bfseries}{\thesection.}{.5em}{\Large}

\titleformat{\chapter}[block]{\normalfont\huge\bfseries}{\thechapter.}{.5em}{\Huge}[{\titlerule[0.8pt]}]
\titlespacing*{\chapter}{0pt}{-19pt}{25pt}
\titleformat{\section}[block]{\normalfont\Large\bfseries}{\thesection.}{.5em}{\Large}

\raggedbottom
\setlength{\emergencystretch}{3em}
% Las tablas técnicas incluyen rutas y entidades largas; estos ajustes evitan ruido de Underfull inocuo.
\hbadness=10000
\vbadness=10000
\newcolumntype{L}[1]{>{\RaggedRight\arraybackslash}p{#1}}
\newcolumntype{Y}{>{\RaggedRight\arraybackslash}X}
\makeatletter
\renewcommand{\es@scroman}[1]{#1}
\makeatother
\pdfsuppresswarningpagegroup=1

% Formato del código fuente con lstlisting
\lstset{
  basicstyle=\ttfamily,
  breaklines=true,
}

% Márgenes
\geometry{
    a4paper,
    margin=2.75cm
}

% Límite de profundidad del índice
\setcounter{tocdepth}{2}

% Indentación de párrafos
\setlength{\parindent}{1cm}

\renewcommand{\lstlistingname}{Extracto de código}
\renewcommand*{\lstlistlistingname}{Índice de extractos de código}

%------------ Caption encima de las tablas
\DeclareCaptionLabelSeparator*{spaced}{\\[1.5ex]}
\captionsetup[table]{textfont=it,format=plain,justification=justified,
  singlelinecheck=false,labelsep=spaced,position=top,skip=8pt}
